\documentclass[11pt,a4paper]{article}
\usepackage[T1]{fontenc}\usepackage[utf8]{inputenc}\usepackage{lmodern}
\usepackage[a4paper,margin=2.2cm]{geometry}
\usepackage{amsmath,amssymb}\usepackage{enumitem}\usepackage{booktabs,longtable,array}
\usepackage[hidelinks]{hyperref}
\setlength{\parindent}{0pt}\setlength{\parskip}{0.4em}
\title{\textbf{OP-matter / A4 portal audit: can the recorded construction
forbid or suppress $\sigma F^2$, $\sigma B^2$, $\sigma W^2$ portals?}\\\large
Recon v1 --- proposal-only}
\author{Per reviewer directive after D3 acceptance}
\date{12 June 2026}
\begin{document}\maketitle

\section*{0. Scope and firewall}
The reviewer's question: is the absence of explicit portal operators
$\sigma F_{AB}F^{AB}$, $\sigma B_{AB}B^{AB}$, $\sigma W_{AB}W^{AB}$ (condition
C-c of G2 / hypothesis H4 of D2) derivable or suppressible from the recorded
matter/gauge construction, or is it an independent completion-level
assumption? Outcome below: \emph{partially derivable} --- tree-level portals
are absent within the recorded induced-gauge architecture, loop-level portals
reduce to the already-enforced charged-sector selection rule, and the
irreducible completion residue is named. One honesty point up front: A4
itself (universal coupling of slow matter to $g^{\rm eff}_{\mu\nu}$) is a
statement about the \emph{metric} channel; portals are non-metric channels, so
A4 alone does not forbid them. The force of this audit comes from the
recorded \emph{induced-gauge} architecture, not from A4 in isolation.
Proposal-only; no preprint edits.

\section{The portal operators and their cost}
The amplitude fluctuation carries mass dimension two ($[\sigma]=[\delta
u]={\rm GeV}^2$), so the leading portals are dimension-four-with-coefficient
operators of the form $(c_P/\Lambda^2)\,\sigma F_{AB}F^{AB}$ with
dimensionless $c_P$ and a completion scale $\Lambda$. On the cosmological
branch the background value $\sigma\!\to\!\delta u_{\rm cos}$ converts the
portal into a $u$-leg of the gauge normalisation,
\[
\gamma_Z^{\rm portal}\;\simeq\;\frac{4\,c_P\,u_0}{\Lambda^2\,Z^{(0)}}\,,
\]
which the D2/G2 thresholds bound at $|\gamma_Z|\lesssim3\times10^{-5}$
(stress level). The implied bound on $c_P$ is scale-sensitive: with the
recorded kernel chain $u_0\sim m_\sigma\phi_0$ and the natural identification
$\Lambda\sim m_\sigma$, one gets $|c_P|\lesssim 3\times10^{-5}\,
Z^{(0)}(m_\sigma/4\phi_0)$ --- an order-of-magnitude, normalisation-dependent
statement, quoted here only to show that an unsuppressed $O(1)$ portal is in
severe tension with the verified anchors under the stated stress-test
assumptions. The audit question is therefore whether the recorded
construction supplies the suppression structurally.

\section{Tree level: no fundamental $F^2$, hence no tree portal}
The recorded bare action (P1--P6) contains only $\Phi$-sector terms; there is
no fundamental gauge field and no fundamental $F_{AB}F^{AB}$. The Maxwell and
electroweak kinetic terms are \emph{emergent/induced} structures: the $U(1)$
Maxwell action is derived in the compact-phase emergence section, the
gravitational and stiffness operators arise as Sakharov-type induced
contributions from integrating out heavy modes (the recorded
induced-structure route), and $(Z_B^{(0)},Z_W^{(0)})$ are defined in the
record as \emph{matching targets at the condensate scale} $\phi_0$, not as
bare couplings. Consequence: \textbf{a tree-level $\sigma F^2$ portal cannot
be written in the recorded theory, because there is no tree-level $F^2$ for
it to dress.} Every gauge-kinetic structure, portal or not, is loop-induced.
This statement is conditional on the recorded induced-gauge reading of the
architecture; a completion that instead introduced fundamental gauge fields
at bare level would re-open the tree channel and would simultaneously abandon
the recorded emergence route.

\section{Loop level: portals reduce to the G2 selection rule}
A loop-induced $\sigma F^2$ coefficient requires charged states running in
the loop whose action carries a $\sigma$-leg --- i.e.\ charged states whose
masses or couplings depend on the amplitude $u$:
\[
c_P^{\rm loop}\;\sim\;\sum_{f\,\in\,\text{charged}}\frac{k_f}{16\pi^2}\,
\frac{\partial\ln m_f^2}{\partial\ln u}\quad(+\;\text{coupling legs}),
\]
which is the \emph{same} object as the threshold-log term in the G2 master
formula. Under the recorded anchoring $m_f=y_fv_2$ with the pinned
$\Phi$-sector evaluation point (and provided the matter-sector matching point
$u_0,\phi_0$ remains distinct from the cosmological $u_{\rm cos}(z)$
excursion), and with the additional flag that the matched Yukawas carry no
$u$-leg ($\gamma_{y_f}=0$), the loop-induced portal vanishes together with
the threshold channel. \textbf{Hence, at loop level, condition C-c is not
independent: it reduces to C-b extended to couplings.} No new threshold is
needed beyond the existing completion-level selection rule
$|\partial\ln m_f^2/\partial\ln u|\lesssim2\times10^{-2}$; a state violating
it would switch on the threshold channel and the portal channel
\emph{together}, at the same quantified $10^{-3}$-per-state level.

\section{The irreducible completion residue}
What the recorded structure does \emph{not} exclude:
\begin{enumerate}[leftmargin=1.9em,itemsep=2pt,label=(R\arabic*)]
\item heavy completion states beyond the recorded spectrum that are
\emph{both} charged and $\sigma$-coupled --- their tree exchange or loops
would induce portals; this class is exactly the class already constrained by
the G2 selection rule, so it adds a name, not a new condition;
\item a microscopic bare-level vertex coupling the magnitude sector directly
to phase-curvature squared (a hand-added $\sigma$--$(\partial\theta)^2$-type
structure outside P1--P6); this would be a deliberate departure from the
recorded bare action and is excluded by bare-action discipline, not by a
theorem.
\end{enumerate}
These two items are the honest residual content of C-c/H4 once the recorded
architecture is taken into account.

\section{Verdict}
\begin{enumerate}[leftmargin=1.9em,itemsep=2pt,label=\textbf{V\arabic*.}]
\item The reviewer's dichotomy is answered: portal absence is
\emph{partially derivable}. Tree-level portals do not exist within the
recorded induced-gauge architecture (no fundamental $F^2$ to dress);
loop-level portals reduce to the charged-sector $u$-decoupling already
enforced by the G2 selection rule, extended to couplings
($\gamma_{y_f}=0$ flag).
\item The independent assumption content of C-c shrinks to the named residue
(R1)--(R2): completion-level charged $\sigma$-coupled mediators, or a
deliberate bare-level magnitude--phase-curvature vertex. Neither exists in
the recorded P1--P6 set.
\item A4's role is supporting, not decisive: it routes the amplitude into
matter through the metric composite, but the portal exclusion itself rests on
the induced-gauge architecture. Stating otherwise would overclaim.
\item D4 candidate (per the review plan, pending review of this audit):
upgrade the OP-CONST block's condition wording from ``absence of
above-threshold $\sigma F^2$ portals'' as an external assumption to
``within the recorded induced-gauge architecture, tree-level portals are
absent and loop-level portals are controlled by the same charged-sector
selection rule; the residual portal assumption is confined to
completion-level states (R1) and bare-action discipline (R2).'' Conditional
throughout; no unconditional-theorem language.
\end{enumerate}
\end{document}
